\scanpage{13}
% typescript upside down in the scan; order follows the typescript
\mspage{PU 78}

Peut-être que \deleted{\unsure{je}} je \handwritten{\inserted{suis en train là}} digresser - car qui me dit que toutes ces ``langue\handwritten{s}'' que je viens d'évoquer à l'instant, font bien partie de cette ``langue-paraboles'', la langue-mère, qui les engloberait toutes\handwritten{?} Je soupçonne pourtant qu'il en est bien ainsi \ldots{} Mais je reviens à\deleted{\unsure{x}} la forme familière, la \handwritten{``}\uline{parabole}\handwritten{''} - \deleted{\unsure{celle}} où la présence d'un \mbox{s\,e\,n\,s}, exprimé par un génial Metteur-en-scène-Prestidigitateur, saute aux yeux (alors même qu'on s'efforcerait à tout prix de l'ignorer). Si pendant longtemps j'ai été réticent \deleted{\unsure{d'}}\handwritten{à} y reconnaître une véritable \mbox{l\,a\,n\,g\,u\,e} (et je sens que cette réticence, maintenant encore, n'est pas entièrement résorbée \ldots), c'est sans doute, justement, parce qu'elle semble si peu apte à ce qui paraît la raison d'être même de toute langue et de tout langage: la communication avec autrui.

Et pourtant, je sais bien qu'il est plus important encore de savoir ``communiquer'' avec \uline{soi-même}.

Quand l'initiative pour ``parler avec moi-même'' vient de ``moi'', à l'état de veille, l'idée\deleted{\unsure{x}} ne me \deleted{serait venue}\handwritten{\inserted{viendrait}} \deleted{de faire appel} (jusqu'à tout récemment encore, du moins) de faire appel à autre chose qu'à la langue-mots, celle donc que ``je'' connais et manie à l'aise (que ce soit en allemand, ou en français). Mais je sais bien aussi que quand l'initiative vient non de moi, mais de ``l'Autre'', ce n'est jamais \deleted{cette} dans cette langue-là qu'\deleted{\unsure{i}}\handwritten{I}l me parle. C'est toujours dans la langue-paraboles, la ``langue des images'' - et quand je prends la peine d'écouter, souvent je sue\handwritten{\inserted{\unsure{r|*}}} sang et eau pour la ``traduire'' dans \deleted{l'autre}\handwritten{\inserted{``la mienne''}} tant bien que mal. Je ne sais si un jour ce ne sera plus la peine \ldots

\deleted{Peut-être est-ce là,\deleted{\unsure{pa}} pour la langue-paraboles, \deleted{mais \unsure{nulle}} (en tant que ``langue'' tout au moins, \deleted{c'est à dire} comme moyen d'une \uline{communication}) sa seule et grande fonction: de} \handwritten{\inserted{Il semblerait que la ``raison d'être'' de la langue-paraboles soit d'être le moyen entre tous pour}} \mbox{n\,o\,u\,s}\quad\mbox{p\,a\,r\,l\,e\,r} à \mbox{n\,o\,u\,s-m\,ê\,m\,e}\quad\mbox{e\,t}\quad\mbox{s\,u\,r}\quad\mbox{n\,o\,u\,s-m\,ê\,m\,e}. C'est la langue \deleted{\unsure{entre toutes}} que choisit notre ``inconscient'', pour ``parler au conscient''\unsure{.|,} \handwritten{L}a langue \handwritten{\inserted{par laquelle \deleted{se font connaître} s'expriment et se font connaître}} les couches créatrices profondes - celles qui ``\uline{savent}'' et qui ``\uline{peuvent}''\handwritten{,}\deleted{\unsure{;}} domicile et demeure de l'\mbox{A\,u\,t\,r\,e}\deleted{\unsure{(}} \handwritten{,} du \mbox{J\,o\,u\,e\,u\,r}, du \mbox{R\,ê\,v\,e\,u\,r}, du \deleted{\unsure{M}} \uline{Metteur-en-scène} ou \uline{Prestidigitateur} ou quelque autre nom qu'on lui donne; l'\mbox{O\,e\,i\,l} qui voit et l'\mbox{O\,r\,e\,i\,l\,l\,e} qui entend et la \mbox{M\,a\,i\,n} qui tient les clefs de tous les \inserted{coffres et tous les} souterrains, et la torche \handwritten{\inserted{aussi}} pour \deleted{\unsure{en}} voir le fond de tous et de chacun \ldots

C'est même étrange que l'Autre jamais ne se lasse de ``nous'' parler, alors que jamais nous ne l'écoutons, autant dire. Souvent pourtant je sens \deleted{qu'Il \unsure{s}} à quel point Il se régale à ses propres jeux - non \deleted{sans doute}\inserted{certes} à les contempler, mais à les inventer et à les monter \inserted{de toutes pièces} avec \handwritten{cette} verve \deleted{\handwritten{\unsure{inventive}}} qui n'est qu'à lui, et sans nul souci, faut\handwritten{-}il croire, s'il y aura un spectateur-auditeur. Il est la \mbox{\handwritten{V}\deleted{\unsure{v}}\,o\,i\,x} de ce qui \uline{est} en nous et qui semble

% archive sheet number: 12
